%% kapa: exam=1 round=30 year=2019 period=1 form=A subject=real_estate subject_ko=부동산학원론 seq=27 number=107
%% source: Q-Net 감정평가사 기출문제 (cst00309) artlSeq=5205205 file=extracted/1008/5205205/2019년 제30회 감정평가사 1차시험 문제지/2019년 제30회 감정평가사 1차시험 문제지/2019년도 제30회 감정평가사 1차 1교시 A형.hwp
%% conversion: kordoc 4.24.51 -> tools/gen1.py (auto); review=proofread-sonnet; flags=none
\begin{question}{107}
부동산시장에 관한 설명으로 옳지 않은 것은? (단, 주어진 조건에 한함)
\begin{choices}
  \item 부동산시장은 단기적으로 수급조절이 쉽지 않기 때문에 가격의 왜곡이 발생할 가능성이 높다.
  \item 부동산의 공급이 탄력적일수록 수요증가에 따른 가격변동의 폭이 크다.
  \item 취득세의 강화는 수급자의 시장진입을 제한하여 시장의 효율성을 저해한다.
  \item 토지이용 규제로 인한 택지공급의 비탄력성은 주택공급의 가격탄력성을 비탄력적으로 하는 요인 중 하나이다.
  \item 주택시장에서 시장균형가격보다 낮은 수준의 가격상한규제는 장기적으로 민간주택 공급량을 감소시킨다.
\end{choices}
\end{question}
\answer{2}
